\BeleskeID{04-s010}
% element: 04-s010:e01 — Množenje RV10 sa r i izdvajanje c−1; razvijene težine i skup cifara
% element: 04-s010:e02 — Naredno množenje razlomljenog ostatka i izdvajanje c−2
% element: 04-s010:e03 — Uslov nultog ostatka i nastavak sa pratećim nulama
% element: 04-s010:e04 — Početno razvijanje svih članova pre množenja sa r
% element: 04-s010:d01

Množenje pomera sve razlomljene težine za jedno mesto ulevo:
\[\begin{aligned}
RV_{10}\cdot r
&=(c_{-1}r^{-1}+c_{-2}r^{-2}+\cdots+c_{-k+1}r^{-k+1}+c_{-k}r^{-k})r\\
&=c_{-1}r^0+c_{-2}r^{-1}+\cdots+c_{-k+1}r^{-k+2}+c_{-k}r^{-k+1}\\
&=c_{-1}+(RV_{10})_{-2}.
\end{aligned}\]
Pošto je početni razlomljeni deo manji od jedinice, proizvod je manji od $r$. Njegov celobrojni deo je zato dozvoljena cifra $c_{-1}$, a preostali razlomljeni deo $(RV_{10})_{-2}$ ponovo je između nule i jedinice.

Za naredni korak množimo samo taj ostatak, ne ceo prethodni proizvod. Sledeći celobrojni deo daje $c_{-2}$. Cifre se ovde dobijaju u konačnom redosledu, od prve iza tačke nadalje. Ako ostatak postane nula, svaki sledeći proizvod ostaje nula i nastaju samo prateće nule.
